\section*{Ζητούμενο 8 - Ανάλυση και Αξιολόγηση της Αλληλεπίδρασης με το LLM}
Η αλληλεπίδραση με το ChatGPT σχετικά με τον κώδικα "GPT from scratch" του Andrej Karpathy ανέδειξε την ικανότητα των σύγχρονων μοντέλων να λειτουργούν ως εργαλεία Explainable AI (XAI), αναλύοντας σύνθετες αρχιτεκτονικές σε πολλαπλά επίπεδα αφαίρεσης.

\subsection*{1. Επεξήγηση Κώδικα και Αρχιτεκτονικής}

Η απόδοση του μοντέλου στην επεξήγηση υπήρξε εξαιρετικά ακριβής, συνδέοντας τον κώδικα PyTorch με τη θεωρία των Transformers.

\begin{itemize}
    \item \textbf{Σημεία-Κλειδιά}: Το LLM διέκρινε ορθά τη ροή δεδομένων από το Character-level Tokenization (χρήση stoi/itos) προς τα Embeddings (Token + Positional).
    \item \textbf{Θεωρητική Σύνδεση}: Ερμήνευσε το Self-Attention ως μηχανισμό «επικοινωνίας» μεταξύ των tokens και το Feed-Forward Network (FFN) ως το στάδιο «επεξεργασίας/σκέψης» της πληροφορίας. Η ικανότητά του να εξηγεί τη σημασία των Residual Connections (x = x + ...) για την αποφυγή του προβλήματος vanishing gradient ευθυγραμμίζεται πλήρως με τη θεωρία των βαθιών δικτύων.
\end{itemize}


\subsection*{2. Αξιολόγηση της Υλοποίησης}
Στο στάδιο της αξιολόγησης, το LLM λειτούργησε ως «έμπειρος κριτικός», αναγνωρίζοντας τα trade-offs μεταξύ παιδαγωγικής σαφήνειας και υπολογιστικής απόδοσης.
\begin{itemize}
    \item \textbf{Σημεία-Κλειδιά}: Εντόπισε ότι ο κώδικας, αν και ιδανικός για εκμάθηση των "first principles", στερείται κλιμακωσιμότητας.
    \item \textbf{Τεχνική Κριτική}: Ανέφερε συγκεκριμένες ελλείψεις έναντι των βιομηχανικών προτύπων (π.χ. GPT-4), όπως η απουσία Flash Attention και η χρήση απλού character-level encoding αντί για Byte-Pair Encoding (BPE), το οποίο θα μείωνε την πολυπλοκότητα $O(n^2)$ του attention mechanism σε μεγάλες ακολουθίες.
\end{itemize}


\subsection*{3. Ανασύνταξη (Refactoring) και Εκσυγχρονισμός}
Το μοντέλο απέδειξε ότι μπορεί να αναβαθμίσει τον κώδικα ενσωματώνοντας σύγχρονες βελτιστοποιήσεις του PyTorch 2.0+.
\begin{itemize}
    \item \textbf{Σημεία-Κλειδιά}: Πρότεινε τη χρήση της scaled\_dot\_product\_attention, η οποία είναι υπολογιστικά ανώτερη από τη χειροκίνητη υλοποίηση του Karpathy.
    \item \textbf{Βελτιώσεις}: Ενίσχυσε το modularity του κώδικα και πρόσθεσε μηχανισμούς σταθερότητας, όπως το Gradient Clipping, έννοιες κρίσιμες για την εκπαίδευση μεγάλων μοντέλων που εξετάζουμε στο εργαστήριο.
\end{itemize}

    
\subsection*{Συμπερασματικές Παρατηρήσεις}
Η συνολική αλληλεπίδραση καταδεικνύει ότι τα LLMs δεν είναι απλώς γεννήτριες κώδικα, αλλά γνωστικοί βοηθοί ικανοί για βαθιά σημασιολογική ανάλυση. Το μοντέλο κατάφερε να γεφυρώσει το χαμηλό επίπεδο υλοποίησης (tensors/matrices) με το υψηλό επίπεδο αρχιτεκτονικής σχεδίασης. Η συνεισφορά του είναι ιδιαίτερα πολύτιμη στον εντοπισμό bottlenecks και στην εφαρμογή βέλτιστων πρακτικών best practices, καθιστώντας το αναπόσπαστο εργαλείο για την κατανόηση της εσωτερικής λειτουργίας των μοντέλων που μελετώνται στην Επεξεργασία Φυσικής Γλώσσας
